\makereporttitle{2023}{octombrie}

\makereportpreamble

\textbf{Rezumat:} Pe parcursul acestei luni am început să îmi completez cunoștințele existente de fizică și să le integrez cu cele de matematică și informatică, ajungând la a înțelege stadiul actual al cercetării în mai multe zone din fizica modernă. În acest scop am participat la cursurile mai multor direcții din cadrul Școlii Doctorale de Fizică și m-am documentat independent din resurse online sau din cărți. 

Totodată, am luat legătura cu diferiți cercetători și profesori din cadrul Facultății de Fizică și din institutele asociate de pe platforma Măgurele și am început să-mi formez legături cu grupurile lor de cercetare.
\vspace{0.5em}

\textbf{Activitate de documentare:}

\begin{itemize}
    \item Am început să parcurg lucrări de cercetare ale dlui.~prof.~Băran despre comportamentul haotic al sistemelor dinamice care apar în fizica cuantică (,,Fingerprints of classical chaotic dynamics in quantum behavior'', ,,Quantifying the folding mechanism in chaotic dynamics'' și ,,Classical and quantal chaos described by a fourth order quadrupole boson hamiltonian'').

    \item Am discutat cu dl.~prof.~Alexandru Nemneș despre direcțiile de cercetare abordate de grupul său, cum ar fi simulările \emph{ab initio} pentru înțelegerea transportului de electroni în semiconductorii bazați pe perovskit sau utilizarea învățării automate pentru a eficientiza simulările bazate pe calcul numeric (dinamică moleculară, prezicerea densității în starea fundamentală).
\end{itemize}
\vspace{0.5em}

\textbf{Cursuri:}
Am participat la mai multe cursuri de la alte direcții de studiu din cadrul Școlii Doctorale:
\begin{itemize}
    \item ,,Diagrame de Fază complexe'', ținut de prof.~dr.~Horia Alexandru;
    \item ,,Tranziții de fază în sisteme macromoleculare'', ținut de prof.~dr.~Valentin Barna;
    \item ,,Atmosfera ca sistem termohidrodinamic'' și ,,Concepte fundamentale în fizica atmosferei'', ținute de prof.~dr.~Sabina Ștefan;
    \item ,,Prelucrarea semnalelor bioelectrice'', ținut de prof.~dr.~Andrei Barborică;
    \item ,,Plasme și metode pentru generarea lor'' și ,,Procese fizice asociate plasmelor și descrierea lor'', ținut de prof.~dr.~Gheorghe Dinescu;
    \item ,,Transportul în nanostructuri bazate pe carbon'', ținut de prof.~dr.~Daniela Dragoman.
\end{itemize}
De asemenea, am auditat cursul ,,Introduction to Symplectic Geometry'' (prelegeri despre formalismul lagrangean și hamiltonian în mecanică, grupuri și algebre Lie), ținut de lect.~dr.~Miron Stanciu (FMI-UB) și am parcurs primele 6 module din cursul online ,,Quantum Mechanics'', oferit de Georgetown University prin intermediul platformei \textit{edX}, ținut de prof.~dr.~Jim Freericks.
\vspace{0.5em}

\textbf{Observațiile conducătorului științific}:

\makereportpostamble{30 octombrie 2023}